\selectlanguage{ngerman}%
\frenchspacing%
\begin{ledgroupsized}[r]{120mm}%
\footnotesize%
\pstart%
\noindent\textbf{Überlieferung:}%
\pend%
\end{ledgroupsized}%
\begin{ledgroupsized}[r]{114mm}%
\footnotesize%
\pstart%
\parindent -6mm%
\makebox[6mm][l]{\textit{L}}%
Aufzeichnung:
LH~I~3,\;8d Bl.\;6.
Ein Blatt 8\textsuperscript{o};
Ränder zum Teil beschnitten.
Eine Seite auf Bl.\;\href{\Scan{LH_01_03_08d_006r}}{6\,r\textsuperscript{o}};
die \href{\Scan{LH_01_03_08d_006v}}{Rückseite} ist unbeschrieben.
\newline%
KK~1 Nr.\;0
\pend%
\end{ledgroupsized}%
\vspace{5mm}%
\begin{ledgroup}%
\footnotesize%
\pstart%
\noindent%
\textbf{Datierungsgründe:}
Die Datierung der Aufzeichnung N.\;\ref{60646} beruht auf Eigenschaften des Textträgers sowie auf inhaltlichen Betrachtungen.
Siehe die ausführliche Begründung auf S.\;\pageref{60646_dat}.%
\pend%
\end{ledgroup}%
\selectlanguage{latin}%
\frenchspacing%
\vspace{8mm}%
\normalsize%
\pstart%
\noindent%
\bracks{%
\href{\Scan{LH_01_03_08d_006r}}%
{6\,r\textsuperscript{o}}%
}
Nova et utilis designatio%
\indexS{designatio}
scribendi Calendarii%
\indexS{calendarium temporum non futurorum, sed praecedentium}
temporum non futurorum%
\indexS{tempus futurum}
sed praecedentium.%
\indexS{tempus praecedens}
Esto aliquis
qui exacte subnotet
quae sit in tractu%
\indexS{tractus quem aliquis habitat}
quem habitat
omnibus totius anni diebus%
\indexS{dies}%
\indexS{annus}
aeris tempestatisque ratio,%
\indexS{ratio aeris tempestatisque}%
\indexS{aer tempestasque}%
\indexS{tempestas}
\edtext{%
quae anni%
}{%
\lemma{quae}%
\Bfootnote{%
\textit{(1)}~tal\angles{i} %
\textit{(2)}~anni%
~\textit{L}%
}}
foecunditas,%
\indexS{foecunditas anni}
salubritas,%
\indexS{salubritas anni}
et quicquid tale a naturalibus causis pendet%
\indexS{quicquid pendens a causis naturalibus}%
\indexS{causa naturalis}
et a vulgaribus Calendariographis%
\indexS{calendariographus vulgaris}
praedici solet.
Calendaria ejusmodi,
fine cuiuslibet anni%
\indexS{annus}
instar Schedularum mortalitatis,%
\indexS{schedula mortalitatis}
edantur;
idque fiat variis in
\edtext{%
locis%
\indexS{locus varius}%
\bracks{,}
quia%
}{%
\lemma{locis}%
\Bfootnote{%
\textit{(1)}~; nam %
\textit{(2)}~quia%
~\textit{L}%
}}
terrarum climatumque constitutio variat.%
\indexS{constitutio terrarum}%
\indexS{constitutio climatum}%
\indexS{terra}%
\indexS{clima}
Hoc si continuo fiat,
tum demum spes est%
\indexS{spes}
\edtext{%
scientiam astrologicam%
\indexS{scientia astrologica}%
\indexS{astrologia}%
}{%
\lemma{scientiam}%
\Bfootnote{%
\textit{(1)}~meteo\indexS{scientia meteorologica} %
\textit{(2)}~astrologicam%
~\textit{L}%
}}
restauratum iri.
Sed et sit
qui subnotet curiose
quae facies coeli,%
\indexS{facies coeli}%
\indexS{coelum}
quae aeris gravitas%
\indexS{gravitas aeris}%
\indexS{aer gravitans}
\edlabel{LH_01_03_08d_006_006r_boy-1}%
\edtext{%
(v.\;Boylei%
\indexN{\textls{Boyle} (Boylius, Boyl), Robert 1627\textendash1691}
Baroscopium),%
\indexS{baroscopium Boylei}%
}{%
\lemma{Boylei Baroscopium}%
\Cfootnote{%
Anspielung auf Boyles\indexN{\textls{Boyle} (Boylius, Boyl), Robert 1627\textendash1691} Verwendung der von ihm verbesserten Vakuumpumpe bei Versuchen über das Gewicht und die Elastizität der Luft.
Siehe insbesondere \textsc{R.\;Boyle},
\textit{New experiments physico-mechanicall}
(Oxford 1660;\cite{00015}
\textit{BW}~I, S.\;141--301);\cite{00017}
Leibniz las vermutlich die lateinische Ausgabe:
\textit{Nova experimenta physico-mechanica}
(Oxford 1661).\cite{00016}%
}}%
\edlabel{LH_01_03_08d_006_006r_boy-2}
quae mox secuta
\edtext{%
tempestas,%
\indexS{tempestas}
quis NB continue ventus spiraverit%
\indexS{ventus spirans}
addito simul situ,%
\indexS{situs}
loci nullis minutissimis circumstantiis omissis%
\indexS{circumstantia loci minutissima}%
\bracks{;}%
}{%
\lemma{tempestas,}%
\Bfootnote{%
\textit{(1)}~nullis minutiss %
\textit{(2)}~quis NB %
\bracks{...} circumstantiis omissis,%
~\textit{L}%
}}
qua ratione spes est%
\indexS{spes}
et Astrologicam scientiam perfectum iri.%
\indexS{scientia astrologica}%
\indexS{astrologia}
\edlabel{LH_01_03_08d_006_006r_fuhr-1}%
\edtext{%
Fuhrmannus%
\indexN{\textls{Fuhrmann} (Fuhrmannus), Stephan 1616\textendash1683}
aliquoties meminit,
se tempestates praecedentis anni annotasse.%
\indexS{tempestas}%
\indexS{annus}%
}{%
\lemma{Fuhrmannus \bracks{...} annotasse}%
\Cfootnote{%
Gemeint ist wohl der Kalendermcher Stephan Fuhrmann (1616--1683) aus Lippstadt,\indexO{Lippstadt} welcher ab 1638 für den baltischen und den norddeutschen Raum jährliche Kalenderhefte veröffentlichte.
Dort bediente er sich sytematischer Himmelsbeobachungen, um seine astrologischen Vorhersagen zu begründen.
Siehe "\href{http://www.presseforschung.uni-bremen.de/dokuwiki/doku.php?id=fuhrmann_stephan}{Fuhrmann, Stephan}", in \textsc{K.-D.\;Herbst}, \textit{Biobibliographisches Handbuch der Kalendermacher von 1550 bis 1750} (online: Fassung vom 8.\;Juli 2024).\cite{02549}%
}}%
\edlabel{LH_01_03_08d_006_006r_fuhr-2}
Haec calendaria non pro vulgo,%
\indexS{calendarium non pro vulgo, sed eruditis}%
\indexS{vulgus}
sed eruditis erunt%
\indexS{eruditi}
et ut scientiae capiant incrementum.%
\indexS{incrementum scientiarum}%
\indexS{scienta capians incrementum}
Sed et multae utiles fieri observationes possent%
\indexS{observatio utilis}
si plures calendariographi%
\indexS{calendariographus ratiocinans de futuris}
qui de futuris ratiocinati sunt,%
\indexS{futura}
conferrentur.
\pend%
\newpage%
\vspace{3.0em}%
\pstart%
\noindent%
\centering%
*     *     *
\pend%
\vspace{3.0em}%
\selectlanguage{german}%
\frenchspacing%
\footnotesize%
\pstart%
\noindent%
\textls{Ausführliche Datierungsgründe:}%
\pagelabel{60646_dat}
Der Träger der vorliegenden undatierten Aufzeichnung N.\;\ref{60646} weist keine Spuren eines Wasserzeichens auf.
Der Schreibduktus und die häufige Verwendung von Schriftkürzeln lassen jedoch eine frühe Abfassung des Textes vermuten, die wohl noch in der Mainzer\indexO{Mainz} Zeit erfolgt sein dürfte.
Dafür sprechen auch folgende Beobachtungen:
\pend%
\pstart%
(1) Die Anspielung -- unter dem Stichwort \textit{Boylei Baroscopium} (\refpassage{LH_01_03_08d_006_006r_boy-1}{LH_01_03_08d_006_006r_boy-2}) -- auf Robert Boyles\indexN{\textls{Boyle} (Boylius, Boyl), Robert 1627\textendash1691} experimentelle Versuche über das Luftgewicht erinnert an die ähnliche Darstellung in der \textit{Hypothesis physica nova}, §\;\bracks{24}--27 (Mainz 1671; \textit{LSB} VI,\;2 N.\;40, S.\;233--235).\cite{00256}
Während der Entstehung dieses Traktats im Laufe des Jahres 1670 soll sich Leibniz ausführlich mit Boyles\indexN{\textls{Boyle} (Boylius, Boyl), Robert 1627\textendash1691} naturwissenschaftlichen Schriften befasst haben (siehe zur Genese der \textit{Hypothesis physica nova} die editorische Erläuterung in \textit{LSB} VI,\;2, S.\;219\,f.).\cite{00256}
Später erweiterte sich das Spektrum der Wissenschaftler, die Leibniz als Koriphäer der experimentellen Erforschung der Lufteigenschaften betrachtete.
Ab dem Sommer 1672 war es vornehmlich Otto von Guericke\indexN{\textls{Guericke} (Gerickius, Gericke), Otto von 1602\textendash1686}, den er nannte, wenn er auf experimentelle Versuche mit der Vakuumpumpe anspielte (vgl. etwa \textit{LSB} VIII,\;1 N.\;39, S.\;296.8--15;\cite{01069} N.\;40, S.\;325.15--326.8).\cite{01265}
\pend%
\pstart%
(2) Die Erwähnung von Stephan Fuhrmanns\indexN{\textls{Fuhrmann} (Fuhrmannus), Stephan 1616\textendash1683} Kalendarien als Beispiele systematischer Wetterbeobachtungen (\refpassage{LH_01_03_08d_006_006r_fuhr-1}{LH_01_03_08d_006_006r_fuhr-2}) ist ebenfalls ein Zeichen dafür, dass die Aufzeichnung N.\;\ref{60646} noch während der Mainzer\indexO{Mainz} Zeit entstanden sein dürfte.
Denn seit seinem Aufenthalt in Paris\indexO{Paris} muss Leibniz gewusst haben, dass auch namhafte Naturforscher wie etwa Edme Mariotte\indexN{\textls{Mariotte}, Edme, Seigneur de Chazeuil ca. 1620\textendash1684} systematische Wetterbeobachtungen durchführten oder wenigstens planten.
Dafür spricht sein Briefaustausch über diesen Themenkreis in seinen ersten Hannoveraner\indexO{Hannover} Jahren.
In seinem Brief vom 7.\;Dezember 1677 erinnert Mariotte\indexN{\textls{Mariotte}, Edme, Seigneur de Chazeuil ca. 1620\textendash1684} daran, dass er seine Wetterbeobachtungen noch weiterführe (\textit{Je continue à faire les observations des vents et de la constitution de l'air}), und lädt Leibniz dazu ein, sich selbst an dieser Unternehmung zu beteiligen, aus der sich bemerkenswerte Entdeckungen ergeben könnten (\textit{LSB} III,\;2 N.\;116, S.\;291.7--15).\cite{02551}
Spätestens zu diesem Zeitpunkt -- wohl aber bereits auch in seinen letzten Pariser Jahren\indexO{Paris} -- hätte Leibniz vermutlich auf Mariottes\indexN{\textls{Mariotte}, Edme, Seigneur de Chazeuil ca. 1620\textendash1684} empirische Forschung verwiesen, um die wissenschaftliche Bedeutung systematischer Wetterbeobachtungen hervorzuheben.
Ab März 1678 hätte er sogar auf die von ihm selbst veranlassten Wetterbeobachtungen verweisen können, wie sie in der Handschrift Ms IV, Nr.\;334 überliefert sind. %
In der Tat ist nach heutigem Wissensstand keine spätere Erwähnung Fuhrmanns\indexN{\textls{Fuhrmann} (Fuhrmannus), Stephan 1616\textendash1683} im Leibniz-Nachlass anzutreffen.
\pend%
\pstart%
(3) Die Aufzeichnung N.\;\ref{60646} weist bemerkenswerte Berührungspunkte mit der auf Mitte 1671 bis Anfang 1672 datierbaren Abhandlung \textit{Directiones ad rem medicam pertinentes} (\textit{LSB} VIII,\;2 N.\;70)\cite{01131} auf.
Dort vertritt Leibniz verwandte Ansichten über die praktische (insbesondere sozialmedizinische) Bedeutung der systematischen Ansammlung empirischer Daten (auch Wetterdaten). %
Die Verwandtschaft wird besonders deutlich in einer Passage der \textit{Directiones ad rem medicam pertinentes}, in der -- ähnlich wie in N.\;\ref{60646} -- die Erforschung "astrologischer" Einflüsse und die Praxis der Kalendermacher miteinander verknüpft werden:
\textit{Man mus achtung geben auf effectus Astrologicos, ob zum exempel wahr was man sagt, daß wenn eine frau durante Eclipsi (solari) gebehre, daß sie und das kind bleibe und was dergleichen traditiones mehr. So mus man auch der Calenderschreiber regeln vom bade, schröpfen aderlaßen, so sie auf den mond und himlische Zeichen applicirt examiniren, nach den von Keplero,\indexN{\textls{Kepler} (Keplerus), Johannes 1571\textendash1630} Campanella,\indexN{\textls{Campanella}, Tommaso 1568\textendash1639} Trew,\indexN{\textls{Trew}, Abdias 1597\textendash1669} und andern gelehrten, vorgeschriebenen wegen.} (\textit{LSB} VIII,\,2 N.\;70, S.\;654.4--9).\cite{01131}
Unmittelbar davor hatte Leibniz dafür plädiert, die Aufnahme von \textit{schedulae mortalitatis} zu systematisieren; dabei solle man \textit{die differentias climatum, terrarum, aeris, etc. genau annotiren laßen, da werden viel admirable dinge heraus kommen} (ebd. S.\;653.25--654.2).\cite{01131}
\pend%
\pstart%
In einer zweiten, verwandten Passage, die sogar eine Anspielung auf die vorliegende Aufzeichnung N.\;\ref{60646} einschließen könnte, veweist Leibniz auf eigene frühere Überlegungen über die Bedeutung einer detallierten "kalendarischen" Aufnahme des täglichen Geschehens:
\textit{Und wenn nach Herren Wrenni,\indexN{\textls{Wren} (Wrennus), Christopher 1632\textendash1723} Hook\indexN{\textls{Hooke} (Hookius, Hook), Robert 1635\textendash1703} und anderer Gedancken eine Historia temporum formirt, oder wie ich offt gedacht, Calender von vergangenen jahren gemacht wurden, so müste man minutim einen ieden annotiren laßen was er fur veranderung dabey an sich empfunden} (ebd. S.\;655.17--20).\cite{01131}
Dass Leibniz bereits in diesen Passagen %
namhafte Naturforscher als Beispiele richtiger wissenschaftlicher Praxis erwähnt, kann als weiteres Zeichen gedeutet werden, dass die \textit{Directiones ad rem medicam pertinentes}\cite{01131} nicht vor der Aufzeichnung N.\;\ref{60646} angefertigt wurden.
\pend%
\pstart%
Diese Beobachtungen bekräftigen insgesamt die eingangs aufgestellte Vermutung, dass die vorliegende Aufzeichnung N.\;\ref{60646} in dem Zeitraum zwischen Entstehung der \textit{Hypothesis physica nova}\cite{00256} und Anfertigung der \textit{Directiones ad rem medicam pertinentes} abgefasst wurde, d.\,h. zwischen Frühjahr 1670 und Anfang 1672.
Weder eine etwas frühere noch eine etwas spätere Datierung sind jedoch auszuschließen.
\pend%
\selectlanguage{latin}%
\frenchspacing%
\normalsize%
